\section{ROM provenance: SHA-256 hashes}\label{supp:romhashes}

We do not redistribute the Atari~2600 game cartridges. They are third-party ROMs
acquired through the \texttt{AutoROM} retrieval procedure, following the practice of the
underlying emulator port~\cite{maier2025vcs}. To let a third party verify that they hold the
exact binaries we ran against, we report a SHA-256 hash for each ROM in the bit-exact
roster. The Data-availability statement points here.

Table~\ref{supp:tab:romhashes} lists all 64 games of the Paper-1 bit-exact roster, so the
hash record covers the six-game core (Table~\ref{supp:tab:coverage}), the 44-game scored set
(\Sectionname\ref{supp:scoredset}), and the wider 54-game labeled breadth set
(Table~\ref{supp:tab:t3coverage}). A dagger ($\dagger$) marks the six core games. Each hash is computed over the raw cartridge image
with \texttt{shasum -a 256}; reproducing it requires only the retrieved binary and a standard
SHA-256 tool. The hashes fingerprint the exact inputs; they are provenance, not a result.

% \hash: typeset a 64-hex-character SHA-256 digest in a narrow column, inserting
% zero-width discretionary breaks between characters so the string wraps cleanly.
% Kept local to this S-file so it stays self-contained (mirrors \jkey in supplement.tex).
\makeatletter
\def\hash#1{{\ttfamily\footnotesize\@hashaux#1\relax}}
\def\@hashaux#1{\ifx\relax#1\else#1\discretionary{}{}{}\expandafter\@hashaux\fi}
\makeatother

\begin{longtable}{@{}p{3.0cm}p{9.0cm}@{}}
\caption{\textbf{SHA-256 hashes of the 64 bit-exact ROMs.} One row per game of the
Paper-1 roster. A dagger ($\dagger$) marks the six core games; the 44 games of the scored set
(\Sectionname\ref{supp:scoredset}) carry interpretability scores, and the remaining rows
document the wider substrate whose ROM binaries are pinned but not scored. Hashes are
computed with \texttt{shasum -a 256} over the retrieved cartridge image and are a provenance
record only.}
\label{supp:tab:romhashes} \\
\toprule
Game & SHA-256 \\
\midrule
\endfirsthead
\multicolumn{2}{@{}l}{\footnotesize Table~\ref{supp:tab:romhashes} (continued)}\\
\toprule
Game & SHA-256 \\
\midrule
\endhead
\bottomrule
\endfoot
\texttt{air\_raid} & \hash{6b5a22ec074a5d870e0693d8717d2f7fbf559ba970995d75882f902bde1d7c9b} \\
\texttt{alien} & \hash{9cd556c7d7f55aeb78fc5add45d74ce889c6e22cba6bb535ef1e8bd6b6c764c1} \\
\texttt{amidar} & \hash{da5b760d3adaeba2cb8fc2367ae96e2b8add4da44cd47199b57ff4d54e682d9f} \\
\texttt{assault} & \hash{d32ca4dd9e849af7b776bc73bf464468839ddab66d7e6efacf877fe6eb841dc2} \\
\texttt{asterix} & \hash{b6161db0530d8445b5ef0e8b700500531e980aa10184e5196355f27a7ba60819} \\
\texttt{asteroids} & \hash{d1a6e808412b4535c7e54fc25414943a9418b1095b4d4969dc33290f7eaa72a4} \\
\texttt{atlantis} & \hash{cef4b7b91ea245f7827604a1c620a5ca7f94288c2c6b54d90a931de8eeb5227d} \\
\texttt{bank\_heist} & \hash{c32daffdb926afaa0b11948c85b0b059e26d29af5f9ce213c439727ccf3f3a8f} \\
\texttt{battle\_zone} & \hash{efd6f639e724a771a8331d7f408609d0dc0854533b3163b602720754267943ee} \\
\texttt{beam\_rider} & \hash{83b30cb52649ddf7a5dd6c3336722421aa435e81e0f85913a1011d81e4e6bb38} \\
\texttt{berzerk} & \hash{cdff0422329de6bdf4ab59176e050fd3fa47d329bc0f065768c4d268838a5fd2} \\
\texttt{bowling} & \hash{dff44a85289f5e9f760254e4806d09038326cdfc42eba3e1cd7c64ed6e0b4dce} \\
\texttt{boxing} & \hash{54760d7e071049fbef09370e136af270bddb5aa2a4fa9be0f404457e49f47377} \\
\texttt{breakout}$\dagger$ & \hash{376323f051c3c373c887fd83abead39d87d844ff283d435f4addbfc1710c6fd5} \\
\texttt{carnival} & \hash{80eab759b6bd4fb59e98d4db1295ad89b128a97d2e20ccb212c744ae642c2f76} \\
\texttt{centipede} & \hash{6def669f54bd886f43026e7fc7f82b9acb92680bb3e3aa68da3767f8797c8cbe} \\
\texttt{chopper\_command} & \hash{055637282252b5378a2b7df573726e85afa91264ffb7243d23fab9518b3c207d} \\
\texttt{crazy\_climber} & \hash{486ad4cc89569d797dc5feff6e95d1d5f5f8e78331a4925ef3126bd2aa1bcf09} \\
\texttt{defender} & \hash{e3c7ed7a073fe8f4b2711bebd624d0b249957daa0ebb93d78823fe36b235ba8c} \\
\texttt{demon\_attack} & \hash{8d72feeb267d23cb5d58bb4511566414884c572d2b4dcba8f6893ad7426442d4} \\
\texttt{double\_dunk} & \hash{c0ea53a91a396638662e2d37261407d6bc9d1b1f660e3101b2098d7ef2bc3190} \\
\texttt{elevator\_action} & \hash{f3b3ad03fd8cf9f348d31c1343580a7a9220d1f06c193165df277e47757a8846} \\
\texttt{enduro} & \hash{ca9bb89755a6d6251fa0ad86983f088a53b4013853963cbeb80dca93466c99fc} \\
\texttt{fishing\_derby} & \hash{84753bc6ecb4dd9f47cd267f1fb15e16b18f7b2b54fbb4ea8207750009171cc1} \\
\texttt{freeway} & \hash{3ef620234b98f22bfd7c06d2467bac9d208e763c9f4309452384d59a741ba893} \\
\texttt{frostbite} & \hash{cbfdad89480def69d922adf16a1d89d55a6cb515929edf74e5bea884c9fb7834} \\
\texttt{gopher} & \hash{b4aff03aeb0fb1c8f4914b5a329f51430f00a2d9350aba24573ad78032dd4697} \\
\texttt{gravitar} & \hash{17a19cd7b5cae2b85f3a51fd10e92e141341f3e4238e9ea641cb97656a3ed0a0} \\
\texttt{hero} & \hash{75bd02ec75458ea4bc2cba9b85bb281c3f9da106123e4600b054f0934e36b81f} \\
\texttt{ice\_hockey} & \hash{06df1fc568d9815cf6ed0d08b27555ee53819c210ee6c786aca9c61b1067b0a2} \\
\texttt{jamesbond} & \hash{995253f0d3a2fac712c6ee3518f19fc5e6a8c800e2fd8dbfe49d45ad367a231f} \\
\texttt{journey\_escape} & \hash{963e2d3da52b6ccaff08c833b2d2cb0329a9b9d5f6b3265fcfd7e0830b914df0} \\
\texttt{kangaroo} & \hash{a2826eaea3a854b63b4ca1d7423f1b3347311e0552c6b1f35065e50b2cbca523} \\
\texttt{krull} & \hash{b40f10d14217093082b00b302bff8c6faa0d64c5221e9be36ea0ceb12477517a} \\
\texttt{kung\_fu\_master} & \hash{3f6501a649ad83e970a25827bd492c56128c36535ae2c96f94bab39b27f939ac} \\
\texttt{montezuma\_revenge} & \hash{69d363a747549817599cace65572700dab9ea33f3a5f91979454fdfbeaa837c0} \\
\texttt{ms\_pacman}$\dagger$ & \hash{dde0b43c5dee7ce1a35ac8e2f625b32a16118512d9be4064115053b972280b14} \\
\texttt{name\_this\_game} & \hash{cc219a7246e8f9454f3bd77345b388fbc2c7343e4ea0c4c864a2eb9b0481882a} \\
\texttt{pacman} & \hash{58e781b472e0583b851b7e56cbfe87bd8294a291e6d39d49d7c999107ed5b918} \\
\texttt{phoenix} & \hash{b01085bc5227a04b974846fe06b4018da84fd3e40ac6d1368f7f6bc07c336dfd} \\
\texttt{pitfall} & \hash{c719b47714d87d17e05de770fab2ab0bdaf6248d411539d6ecdeddf7de3ae77b} \\
\texttt{pong}$\dagger$ & \hash{41623e3c26148f96cbe77d56512845d224a63c8fe5bf90e38527155e39ec96d3} \\
\texttt{pooyan} & \hash{ac9abbfc272deb0305451afa3d5cd58849684691cfa25cede01251fab36c131b} \\
\texttt{private\_eye} & \hash{c8f9ce1b2e804b6778dbc20da0c935b2af7d7e998ad372d44df66f16eb4b98ce} \\
\texttt{qbert}$\dagger$ & \hash{3257221832a7607696b06b9985103c12a4c24c61b0c655e8f52aed8e38194b76} \\
\texttt{riverraid} & \hash{4c6842b8af64fc75e599b773821b3e6bee8e5028f5d63414821b0d63f0b2dfe9} \\
\texttt{road\_runner} & \hash{4b6c2e68d6933b8bc79983f6a3bb11bda695b5dfbeeb8fd63d84406ba3f27b5a} \\
\texttt{robotank} & \hash{c186409481e688401c8a7a9d5f0387ed448d18e297e287fd7c409752fa9a2a0d} \\
\texttt{seaquest}$\dagger$ & \hash{fbc29f4678f69ac27fe46da298c8c22f98cba2a3e5491e6b308fd9bf68d2ee43} \\
\texttt{skiing} & \hash{804eff5f0196375b33614aa18cb7f2f883681ced1654e5e8ca702204756a19bb} \\
\texttt{solaris} & \hash{0afa36e5f4d8d77e38716673c96cda3ad11d421283f5a5ae77ec7a01e7718a4e} \\
\texttt{space\_invaders}$\dagger$ & \hash{7224b17462b992d67f4e06a3c85f269c9822b06df6015bf038b55f384ced0301} \\
\texttt{star\_gunner} & \hash{3ad501a39280de82bd604c6d8ae915fe93c1e6ebfd97685ab097c5dbef9822c3} \\
\texttt{surround} & \hash{c6864f5c9f43f0590226b4f2943d00fb8a90a2c6f799963de75016d5c07affc3} \\
\texttt{tennis} & \hash{5a2052f020bfb19f5fa4db8a3149287a875bc4025059dc66b747451fb7fe4909} \\
\texttt{time\_pilot} & \hash{8ea97e335a2fb122faa591a86e695de1eddb88e9370d3ead6e82d308f6e81c7d} \\
\texttt{tutankham} & \hash{403397fe85838ebc94bf063e63756b7656b9c962b9c6392b991d37ef4ad4a3d4} \\
\texttt{up\_n\_down} & \hash{74890a43e59134ea88e0a3aa6cae91615299cc41d74a464e0e31070d5d36062c} \\
\texttt{venture} & \hash{1870b0dd6fb1de45e44e1ad114c8f558104351cfb11495dc4fa5c7b0fc58f455} \\
\texttt{video\_pinball} & \hash{cef82c39bbb035f1ab5b900548de410371753ff717f3bf69a2e314cbd55fe937} \\
\texttt{videochess} & \hash{d44f2a115393b0a6cbcd5398d62369e652f3670a9ee585883d183e740edad0f7} \\
\texttt{wizard\_of\_wor} & \hash{d69ee89e65c34cd3152d98d778de913dd7f59e7309f2bcbd3ebead585065e747} \\
\texttt{yars\_revenge} & \hash{ff777c8d4ea06569156fcb2f5ba0e00eebaf2c2c509e5829109e5aefed199c66} \\
\texttt{zaxxon} & \hash{e1c323ce00cccdd4bc737fd44abcef35b24fee04fd90bf9c9374e18b3b2ae6bf} \\

\end{longtable}
